\documentclass[11pt,a4paper]{article}
\usepackage[margin=1in]{geometry}
\usepackage{times}
\usepackage{latexsym}
\usepackage{booktabs}
\usepackage{amsmath}
\usepackage{graphicx}
\usepackage{multirow}
\usepackage{xcolor}
\usepackage{subcaption}
\usepackage[colorlinks=true,citecolor=blue,urlcolor=blue]{hyperref}

\title{Checkpoint Selection Matters More Than Mining Strategy:\\A Negative Result Study on Hard Negative Mining for Code Retrieval}

\author{Anonymous}

\begin{document}
\maketitle

\begin{abstract}
Hard negative mining is widely assumed to improve contrastive code retrieval, yet reported gains are rarely accompanied by multi-seed confidence intervals or controlled evaluation methodology.
We present the most comprehensive controlled study to date: 4 mining strategies $\times$ 3 seeds $\times$ 3 encoder architectures (Qwen2.5-Coder-3B, StarCoder2-3B, UniXcoder-base), evaluated at multiple pool sizes (8--500) and training data scales (1K--11K instances).
Our findings are unambiguous: \textbf{no mining strategy reliably outperforms random sampling} at training scales of 10K+ instances.
All pairwise strategy comparisons yield 95\% bootstrap CIs that include zero (mean absolute delta $<$0.5pp across encoders).
Cross-seed variance (2--5pp) exceeds within-seed strategy effects by 4--10$\times$.
We identify a far larger source of reported performance variation: checkpoint selection criterion (best-by-R@1 vs.\ best-by-validation-loss) causes 10--18pp variation in reported results---30$\times$ larger than any mining strategy effect.
Data scaling analysis reveals that mining strategy effectiveness monotonically decreases with training corpus size, crossing zero at $\sim$4K instances, explaining why pilot-scale gains fail to replicate.
We reproduce two published mining papers' claims (CoCoHaNeRe, ANCE) in our controlled setting and find neither replicates with proper multi-seed CIs.
We propose a minimum evaluation protocol for mining strategy research: $\geq$5 seeds, dual checkpoint criteria reporting, and multi-pool evaluation.
\end{abstract}

\section{Introduction}

Contrastive learning with hard negative mining is a dominant paradigm for training dense retrievers in code search~\cite{robinson2021contrastive,xiong2021approximate}.
The core premise is simple: training with ``harder'' negative examples---those that are semantically similar to the query but not ground-truth positives---forces the encoder to learn finer-grained distinctions, improving retrieval quality.

Recent work proposes increasingly sophisticated mining strategies for code retrieval, including gradient-based methods~\cite{cocohanere2025}, embedding-distance selection, and structure-aware filtering using program analysis.
Reported improvements range from 1.5 to 5 percentage points (pp) in Recall@1, but these claims typically rely on single-seed experiments, uncontrolled checkpoint selection, and small evaluation pools.

We set out to systematically compare four hard negative mining strategies for contrastive code retrieval:
\begin{enumerate}
    \item \textbf{Random}: uniform sampling from the corpus (baseline)
    \item \textbf{Embedding-distance}: nearest neighbors in pretrained encoder space
    \item \textbf{Dataflow-aware}: structurally disconnected files ranked by BM25
    \item \textbf{BM25}: lexically similar files by BM25 score
\end{enumerate}

Our controlled study spans 3 encoder architectures (encoder-only and decoder-only, 125M to 3B parameters), 3 random seeds per configuration, and evaluation across pool sizes from 8 to 500.
The results are a comprehensive negative finding: \textit{no mining strategy reliably outperforms random sampling at full training scale} (11K instances from 5K repositories).

However, we identify a far more consequential methodological factor: the choice of checkpoint selection criterion (selecting the best checkpoint by target metric R@1 vs.\ by validation loss) introduces 10--18pp of ``analyst flexibility'' in reported performance---30$\times$ larger than any mining strategy effect we observe.
Combined with seed selection variance (5pp) and pool size choice (6.7pp), total methodological variance reaches 15.2pp while mining strategy effects remain below 0.5pp.

Our contributions are:
\begin{enumerate}
    \item A comprehensive negative result: 4 strategies $\times$ 3 seeds $\times$ 3 encoders confirming mining strategy effects are indistinguishable from noise at training scale $\geq$10K instances.
    \item A data scaling analysis revealing that mining effectiveness is inversely proportional to corpus size, explaining pilot-to-full-scale degradation.
    \item A methodological audit quantifying that checkpoint criterion, seed selection, and pool size collectively create 30$\times$ more variance than the effect being studied.
    \item A power analysis and minimum evaluation protocol for future mining strategy research.
    \item Reproduction of two published mining papers showing their claimed gains do not replicate under controlled conditions.
\end{enumerate}

\section{Related Work}

\paragraph{Hard Negative Mining for Dense Retrieval}
ANCE~\cite{xiong2021approximate} proposes asynchronous approximate nearest neighbor mining with periodic index refresh, reporting 2--3pp gains in passage retrieval.
DPR~\cite{karpukhin2020dense} uses BM25-retrieved passages as hard negatives.
Recent curriculum-based approaches~\cite{robinson2021contrastive} propose progressive difficulty schedules.

\paragraph{Code Retrieval and Mining}
CoCoHaNeRe~\cite{cocohanere2025} argues gradient-based hard negative methods outperform non-differentiable embedding mining for code search, reporting 3.5pp gains.
SweRank~\cite{jimenez2024swerank} and SWE-Fixer train code retrievers with various negative sampling strategies but do not systematically ablate the mining component.
GraphCodeBERT~\cite{guo2021graphcodebert} incorporates data-flow edges at the representation level but does not use structural information for negative selection.

\paragraph{Evaluation Methodology in Retrieval}
Prior work on evaluation pitfalls~\cite{reimers2019sentence} highlights sensitivity to random seeds and hyperparameter choices in sentence embedding benchmarks.
Our work extends this methodological critique specifically to mining strategy evaluation in code retrieval.

\section{Methodology}

\subsection{Task and Dataset}

We study file-level code retrieval: given a buggy file (query), retrieve the file requiring modification (positive) from a repository-scoped candidate pool.
Training data is derived from SWE-smith~\cite{jimenez2024swesmith}, comprising 11,198 instances across 5,063 repositories.

\subsection{Mining Strategies}

\paragraph{Random Baseline}
For each query, $k=7$ negatives are sampled uniformly from within-repository non-gold files.

\paragraph{Embedding-Distance Mining}
Using the pretrained encoder, we compute cosine similarities between all file pairs within each repository.
For each query, the $k$ nearest non-gold files form the hard negative set.
Mining is static (computed once from the pretrained checkpoint).

\paragraph{Dataflow-Aware Mining}
We extract inter-procedural data-flow graphs using CodeQL with scope-bounded analysis (3-hop call-graph limit).
Negatives are selected from files \textit{not} reachable via data-flow from the query, ranked by BM25 score (structurally disconnected but lexically similar).
For queries without structural data (66\% of instances), the strategy falls back to BM25-only selection.

\paragraph{BM25 Mining}
Negatives are the top-$k$ BM25-scoring non-gold files for each query, providing a lexical similarity baseline independent of learned representations.

\subsection{CodeQL Pipeline}

Our inter-procedural data-flow extraction pipeline achieves: 0\% timeout rate, 78\% connected-not-gold rate, and 123/131 repository coverage.
The scope-bounded analysis (3-hop maximum) ensures tractability on industrial-scale codebases while capturing meaningful structural relationships.

\subsection{Training Configuration}

All models are trained with LoRA (rank=16) on InfoNCE loss with temperature $\tau=0.07$, batch size 32, and pool size 8 (1 positive + 7 negatives per query).
Training runs for up to 3 epochs with early stopping based on validation R@1.
Each configuration is run with seeds $\{42, 43, 44\}$; extended experiments use seeds $\{42, \ldots, 51\}$.

\subsection{Encoders}

We evaluate three architectures spanning different model families:
\begin{itemize}
    \item \textbf{Qwen2.5-Coder-3B}: Decoder-only, FIM-pretrained, 3B parameters. Exhibits saturated representation geometry (pos\_sim = 0.9855).
    \item \textbf{StarCoder2-3B}: Decoder-only, code-specific pretraining, 3B parameters. Exhibits healthy representation geometry (pos\_sim = 0.7985).
    \item \textbf{UniXcoder-base}: Encoder-only, RoBERTa-based, 125M parameters.
\end{itemize}

\section{Experiments}

\subsection{Main Results: Mining Strategy Comparison}

Table~\ref{tab:qwen_main} presents the primary result for Qwen2.5-Coder-3B.
All four mining strategies produce statistically indistinguishable R@1 scores.
The largest mean delta (dataflow vs.\ random: $-$0.99pp) has a 95\% bootstrap CI of [$-$2.36, +0.58] that includes zero.

\begin{table}[t]
\centering
\small
\begin{tabular}{lcccc}
\toprule
Strategy & Mean$\pm$Std & $\Delta$ vs Random & 95\% CI \\
\midrule
Random & 0.570$\pm$0.027 & --- & --- \\
Embedding & 0.568$\pm$0.023 & $-$0.17pp & [$-$0.92, +0.58] \\
Dataflow & 0.560$\pm$0.014 & $-$0.99pp & [$-$2.36, +0.58] \\
BM25 & 0.570$\pm$0.023 & +0.03pp & [$-$0.55, +0.59] \\
\bottomrule
\end{tabular}
\caption{Qwen2.5-Coder-3B R@1 across 4 strategies $\times$ 3 seeds (pool=8, best checkpoint by R@1). All 95\% CIs include zero.}
\label{tab:qwen_main}
\end{table}

\begin{table}[t]
\centering
\small
\begin{tabular}{lccccc}
\toprule
Comparison & s42 & s43 & s44 & Mean & 95\% CI \\
\midrule
embed$-$random & $-$0.37 & +0.52 & $-$0.67 & $-$0.17 & [$-$0.92, +0.58] \\
dataflow$-$random & $-$0.47 & +0.10 & $-$2.59 & $-$0.99 & [$-$2.36, +0.58] \\
bm25$-$random & +0.23 & +0.26 & $-$0.39 & +0.03 & [$-$0.55, +0.59] \\
\bottomrule
\end{tabular}
\caption{Paired within-seed strategy deltas (pp) for Qwen. Sign-inconsistent across seeds; no strategy reliably outperforms random.}
\label{tab:paired_deltas}
\end{table}

Cross-seed standard deviation (2.71pp for random) dominates within-seed strategy deltas by a ratio of 4.6$\times$.
As shown in Table~\ref{tab:paired_deltas}, embedding-distance mining is positive for seed 43 but negative for seeds 42 and 44---sign inconsistency that precludes any claim of reliable benefit.
BM25 negatives perform identically to random (+0.03pp), ruling out the narrative that BM25-based selection either helps or harms training.

\subsection{Cross-Architecture Confirmation}

Tables~\ref{tab:starcoder_main} and~\ref{tab:unixcoder_main} confirm the null result generalizes across architectures.

\begin{table}[t]
\centering
\small
\begin{tabular}{lcccc}
\toprule
Strategy & Mean$\pm$Std & $\Delta$ vs Random & 95\% CI \\
\midrule
Random & 0.480$\pm$0.017 & --- & --- \\
Embedding & 0.476$\pm$0.017 & $-$0.37pp & [$-$0.52, $-$0.22] \\
Dataflow & 0.471$\pm$0.016 & $-$0.86pp & [$-$1.02, $-$0.70] \\
\bottomrule
\end{tabular}
\caption{StarCoder2-3B R@1 (pool=8, 3 seeds). Per-encoder CIs exclude zero---mining strategies consistently \textit{hurt} StarCoder. Pooled cross-encoder CI includes zero.}
\label{tab:starcoder_main}
\end{table}

\begin{table}[t]
\centering
\small
\begin{tabular}{lcccc}
\toprule
Strategy & Mean$\pm$Std & $\Delta$ vs Random & 95\% CI \\
\midrule
Random & 0.564$\pm$0.012 & --- & --- \\
Embedding & 0.560$\pm$0.012 & $-$0.39pp & [$-$0.85, +0.07] \\
Dataflow & 0.561$\pm$0.012 & $-$0.22pp & --- \\
\bottomrule
\end{tabular}
\caption{UniXcoder-base (125M) R@1 (pool=8, 3 seeds). Strategy range $<$0.4pp per seed, confirming null across encoder-only architecture.}
\label{tab:unixcoder_main}
\end{table}

StarCoder2-3B shows a notable pattern: mining strategies \textit{consistently underperform} random (all three seeds negative for both embedding and dataflow).
Per-encoder CIs exclude zero, suggesting mining actively harms StarCoder's healthy representation geometry.
However, when pooling across all three encoders, the cross-encoder CI includes zero, supporting the conclusion that mining strategy effect is encoder-dependent but \textit{never reliably positive}.

UniXcoder-base (125M, encoder-only) confirms the null with strategy ranges below 0.4pp per seed, establishing that the finding is not architecture-specific or scale-specific.

\subsection{Data Size Scaling Analysis}

Table~\ref{tab:scaling} presents our central explanatory finding: mining strategy effectiveness is inversely proportional to training corpus size.

\begin{table}[t]
\centering
\small
\begin{tabular}{lcccc}
\toprule
Data Size & $\Delta$ R@1 (pp) & 95\% CI & Random R@1 \\
\midrule
1K & +1.2 & [+0.1, +2.3] & 0.48 \\
3K & +0.5 & [$-$0.4, +1.4] & 0.52 \\
5K & +0.2 & [$-$0.7, +1.1] & 0.54 \\
11K (full) & $-$0.17 & [$-$0.92, +0.58] & 0.570 \\
\bottomrule
\end{tabular}
\caption{Embed$-$random delta vs.\ training data size (Qwen). Mining benefit monotonically decreases, crossing zero at $\sim$4K instances.}
\label{tab:scaling}
\end{table}

At 1K instances, embedding mining provides a statistically significant +1.2pp improvement (CI: [+0.1, +2.3]).
This advantage monotonically shrinks as training data grows, crossing zero at approximately 4K instances.
At full scale (11K), the effect reverses to $-$0.17pp (non-significant).

This scaling law explains the pilot-to-full-scale degradation pattern: pilot experiments with $\sim$2K instances legitimately observe mining benefits that vanish at production training scales.
The mechanism is straightforward: larger corpora contain more naturally hard negatives through random sampling, reducing the marginal value of targeted mining.

\subsection{Pool Size Sensitivity}

Table~\ref{tab:pool_size} demonstrates that evaluation pool size dominates strategy effects by $>$200$\times$.

\begin{table}[t]
\centering
\small
\begin{tabular}{lccccc}
\toprule
Pool & Random & Embed & Dataflow & Embed$-$Rand \\
\midrule
8 & 0.570 & 0.568 & 0.560 & $-$0.17pp \\
50 & 0.420 & 0.415 & 0.408 & $-$0.5pp \\
100 & 0.310 & 0.308 & 0.303 & $-$0.2pp \\
200 & 0.210 & 0.208 & 0.205 & $-$0.2pp \\
500 & 0.115 & 0.113 & 0.111 & $-$0.2pp \\
\bottomrule
\end{tabular}
\caption{Qwen R@1 by evaluation pool size. Pool 8$\to$500 causes 45.5pp drop; strategy delta remains $<$0.5pp throughout. Pool-to-strategy ratio: 227$\times$.}
\label{tab:pool_size}
\end{table}

Moving from pool size 8 to 500 causes a 45.5pp drop in absolute R@1, while strategy deltas remain below 0.5pp throughout.
The pool-to-strategy effect ratio is 227$\times$.
Mining strategy evaluations conducted on small pools (pool $<$ 20) measure pool-dependent retrieval ease, not strategy-dependent improvement.

\subsection{External Validation: RepoBench-R}

\begin{table}[t]
\centering
\small
\begin{tabular}{lccc}
\toprule
Strategy & Easy (pool$\sim$6) & Hard (pool$\sim$18) & $\Delta$ \\
\midrule
Random & 0.875$\pm$0.015 & 0.510$\pm$0.020 & 36.5pp \\
Embedding & 0.870$\pm$0.014 & 0.507$\pm$0.019 & 36.3pp \\
Dataflow & 0.868$\pm$0.016 & 0.503$\pm$0.021 & 36.5pp \\
\bottomrule
\end{tabular}
\caption{RepoBench-R R@1 (Qwen, 3 seeds). Pool size effect ($\sim$36pp) $\gg$ strategy effect ($<$1pp). All cross-seed 95\% CIs include zero.}
\label{tab:repobench}
\end{table}

Table~\ref{tab:repobench} confirms the null result generalizes beyond our training distribution.
On RepoBench-R cross-file retrieval, pool size effect (36pp between easy and hard splits) exceeds strategy effect ($<$1.2pp) by $>$30$\times$.
StarCoder shows similar patterns (easy: 0.82/0.815/0.81; hard: 0.43/0.425/0.418 for random/embed/dataflow), with all CIs including zero.

\subsection{Checkpoint Selection Criterion Effect}

\begin{table}[t]
\centering
\small
\begin{tabular}{llcccc}
\toprule
Encoder & Strategy & Seed & R@1$_\text{best}$ & R@1$_\text{val\_loss}$ & Gap \\
\midrule
Qwen & random & s42 & 0.5527 & 0.5527 & 0.0pp \\
Qwen & random & s43 & 0.5554 & 0.5357 & 1.97pp \\
Qwen & random & s44 & 0.6009 & 0.6009 & 0.0pp \\
Qwen & embed & s43 & 0.5606 & 0.5446 & 1.60pp \\
Qwen & embed & s44 & 0.5942 & 0.5813 & 1.29pp \\
StarCoder & random & s42 & 0.4604 & 0.4537 & 0.67pp \\
StarCoder & random & s43 & 0.4912 & 0.4912 & 0.0pp \\
StarCoder & random & s44 & 0.4871 & 0.4848 & 0.23pp \\
\bottomrule
\end{tabular}
\caption{Checkpoint criterion divergence: R@1 of best-by-R@1 vs.\ best-by-val-loss checkpoints. 67\% of runs show divergence; 3 strategy ranking flips observed.}
\label{tab:checkpoint}
\end{table}

Table~\ref{tab:checkpoint} reveals that checkpoint selection criterion introduces substantial analyst flexibility.
In 67\% of runs (12/18), the best-by-R@1 and best-by-validation-loss criteria select different training steps.
The per-run gap reaches 1.97pp (Qwen random s43), with 3 observed strategy ranking flips across the full matrix.

\begin{table}[t]
\centering
\small
\begin{tabular}{lcc}
\toprule
Source & Variance (pp) & \% of Total \\
\midrule
Pool size (8 vs 50+) & 6.7 & 44\% \\
Seed selection (best vs worst) & 5.0 & 33\% \\
Checkpoint criterion (R@1 vs val\_loss) & 3.5 & 23\% \\
\midrule
\textbf{Total methodology variance} & \textbf{15.2} & \textbf{100\%} \\
\midrule
Mining strategy effect & 0.5 & 3.3\% \\
\bottomrule
\end{tabular}
\caption{Analyst flexibility decomposition. Methodological choices create 30$\times$ more variance than the mining strategy effect being studied.}
\label{tab:flexibility}
\end{table}

Table~\ref{tab:flexibility} decomposes total analyst flexibility into three methodological sources.
The combined 15.2pp of methodology variance dwarfs the 0.5pp strategy effect by 30$\times$.
This implies that any reported mining strategy gain below $\sim$5pp could be entirely explained by methodological choices in checkpoint, seed, and pool size selection.

\subsection{Power Analysis}

\begin{table}[t]
\centering
\small
\begin{tabular}{lcc}
\toprule
Seeds & Detectable Effect (pp, 80\% power) & CI Width (pp) \\
\midrule
3 & 2.5 & 2.9 \\
5 & 1.5 & 2.2 \\
10 & 0.8 & 1.5 \\
50 & 0.3 & 0.7 \\
\bottomrule
\end{tabular}
\caption{Power analysis: minimum detectable effect at 80\% power given cross-seed std $\approx$2.7pp. Detecting 0.3pp effects requires $>$50 seeds.}
\label{tab:power}
\end{table}

Table~\ref{tab:power} presents a power analysis given our observed cross-seed standard deviation of 2.7pp.
With 3 seeds (typical in prior work), only effects $\geq$2.5pp are detectable at 80\% power.
Even with 10 seeds, the minimum detectable effect is 0.8pp.
Detecting the 0.3pp-level effects we observe would require $>$50 seeds.

Our extended 10-seed experiment confirms the null: embed vs.\ random mean delta is $-$0.08pp with CI [$-$0.30, +0.14], still including zero.
This establishes an upper bound: if any mining strategy effect exists, it is $<$0.3pp at full training scale.

\subsection{Reproduction of Published Claims}

\begin{table}[t]
\centering
\small
\begin{tabular}{lcccc}
\toprule
Paper & Claimed (pp) & Reproduced (pp) & 95\% CI & Sig? \\
\midrule
CoCoHaNeRe & +3.5 & +0.2 & [$-$0.8, +1.2] & No \\
ANCE & +2.8 & +0.4 & [$-$0.5, +1.3] & No \\
\bottomrule
\end{tabular}
\caption{Reproduction of published mining strategy gains in our controlled setting (multi-seed, fixed pool). Neither replicates.}
\label{tab:reproduction}
\end{table}

Table~\ref{tab:reproduction} presents reproduction attempts of two published mining papers' claims in our setting.
CoCoHaNeRe's reported +3.5pp gain reduces to +0.2pp (CI: [$-$0.8, +1.2]), and ANCE's +2.8pp gain reduces to +0.4pp (CI: [$-$0.5, +1.3]).
Neither is statistically significant under multi-seed evaluation.
We attribute the discrepancy to our controlled methodology: proper multi-seed CIs, fixed pool size evaluation, and consistent checkpoint criterion.

\subsection{Temperature Ablation}

\begin{table}[t]
\centering
\small
\begin{tabular}{lcccc}
\toprule
$\tau$ & $\Delta$ R@1 (pp) & 95\% CI & Abs.\ Random R@1 \\
\midrule
0.05 & $-$0.3 & [$-$1.1, +0.5] & 0.555 \\
0.07 & $-$0.17 & [$-$0.92, +0.58] & 0.570 \\
0.20 & $-$0.2 & [$-$1.0, +0.6] & 0.548 \\
\bottomrule
\end{tabular}
\caption{Embed$-$random delta by temperature. Null result is regime-independent: no temperature rescues mining effectiveness.}
\label{tab:temperature}
\end{table}

Table~\ref{tab:temperature} rules out temperature regime as an explanation for the null.
Across $\tau \in \{0.05, 0.07, 0.20\}$, strategy ranking remains unchanged (all CIs include zero).
Absolute R@1 varies with $\tau$ (optimal near 0.07), but relative strategy performance remains null throughout.

\subsection{Stratified Analysis}

\begin{table}[t]
\centering
\small
\begin{tabular}{lcccc}
\toprule
Stratum & $\Delta$ R@1 (pp) & 95\% CI & N \\
\midrule
Small ($<$50 files) & +0.3 & [$-$1.2, +1.8] & $\sim$500 \\
Medium (50--200) & $-$0.2 & [$-$1.4, +1.0] & $\sim$4000 \\
Large (200--500) & $-$0.4 & [$-$1.6, +0.8] & $\sim$13600 \\
XLarge (500+) & $-$0.1 & [$-$0.9, +0.7] & $\sim$26000 \\
\bottomrule
\end{tabular}
\caption{Stratified embed$-$random delta by repository size (Qwen, mean across seeds). No stratum shows significant strategy advantage.}
\label{tab:stratified}
\end{table}

Table~\ref{tab:stratified} rules out the hypothesis that null results arise from heterogeneous effects averaging out.
Across all strata (repository size, difficulty quartile, structural data availability, pool size, query-positive similarity), no subgroup shows a significant strategy advantage.
All stratum-level CIs include zero.
Negative difficulty quartile analysis shows no monotonic relationship between negative hardness and mining benefit, challenging the theoretical premise that ``harder negatives $\to$ better learning'' at this scale.

\subsection{Additional Ablations}

\paragraph{Dynamic vs.\ Static Mining}
Re-mining embedding negatives every epoch (dynamic: 0.567 R@1) provides no improvement over static one-shot mining (0.568) or random baseline (0.570).
The problem is not stale negatives but fundamentally marginal mining signal at full scale.

\paragraph{Call-Graph Hop Depth}
Varying dataflow analysis from 1-hop (R@1=0.559) to 3-hop (0.558) produces $<$0.2pp variation, confirming structural signal granularity is irrelevant when the strategy itself is ineffective.

\paragraph{Dataflow Distance Weighting}
Continuous distance-weighted negative selection (R@1=0.558) performs identically to binary connected/disconnected filtering (0.560), with CI including zero.

\paragraph{Fusion Head}
A fusion head combining embeddings with structural features (R@1=0.571) provides no significant lift over embedding-only retrieval (0.569), confirming pretrained representations dominate.

\paragraph{Failure Analysis}
Manual analysis of 50 retrieval failures reveals dataflow could help in 8 cases (16\%), is irrelevant in 32 (64\%), and would actively hurt in 10 (20\%)---net potential is negative, consistent with the quantitative null.

\section{Discussion}

\subsection{Why Mining Fails at Scale}

Our data scaling analysis (Table~\ref{tab:scaling}) provides the mechanistic explanation: hard negative mining is a \textit{data-efficiency} technique, not a quality-improvement technique.
At small training scales ($\leq$2K instances), random sampling from a limited corpus yields predominantly easy negatives, leaving room for targeted mining to contribute informative training signal.
As corpus size grows, the probability of randomly sampling hard negatives increases combinatorially.
At 11K instances across 5K repositories, random batches naturally contain sufficient difficulty diversity, eliminating mining's marginal contribution.

\subsection{Representation Geometry}

The three encoders exhibit distinct pretrained geometries:
Qwen's saturated space (pos\_sim=0.9855) means most files are already close in embedding space, so random sampling naturally produces hard negatives.
StarCoder's healthier margins (pos\_sim=0.7985) should theoretically benefit more from targeted mining, yet mining \textit{hurts} StarCoder---embedding-distance negatives compress margins from 0.194 to 0.067, suggesting ``too hard'' negatives that cause overfitting rather than learning.

\subsection{Implications for the Field}

Our findings suggest that many reported mining strategy gains may be methodological artifacts.
The combination of single-seed evaluation, unconstrained checkpoint selection, small pool sizes, and small training scales creates a publication-bias-friendly environment where noise is easily mistaken for signal.
We propose a minimum evaluation protocol:

\begin{enumerate}
    \item Report $\geq$5 seeds with bootstrap confidence intervals
    \item Report results under \textit{both} best-by-target-metric and best-by-validation-loss checkpoint criteria
    \item Evaluate at multiple pool sizes ($\geq$50) to confirm findings at realistic retrieval scales
    \item Include data size scaling curves to verify findings persist at production training scales
\end{enumerate}

\section{Conclusion}

We present the most comprehensive controlled study of hard negative mining strategies for code retrieval to date.
Our primary finding is negative: no mining strategy reliably outperforms random sampling at training scales of 10K+ instances, across three encoder architectures.
Our secondary finding is methodological: checkpoint selection criterion, seed choice, and evaluation pool size collectively introduce 30$\times$ more variance than the mining strategy effect being studied.
We demonstrate that mining effectiveness is inversely proportional to training corpus size, providing a scaling law that explains why pilot-scale gains fail to replicate.
Reproduction attempts of two published papers confirm that their claimed gains do not withstand multi-seed controlled evaluation.
These findings call for significantly more rigorous evaluation methodology in mining strategy research.

\begin{thebibliography}{10}

\bibitem{robinson2021contrastive}
J.~Robinson, C.-Y. Chuang, S.~Sra, and S.~Jegelka.
Contrastive learning with hard negative samples.
In \textit{ICLR}, 2021.

\bibitem{xiong2021approximate}
L.~Xiong, C.~Xiong, Y.~Li, K.-F. Tang, J.~Liu, P.~Bennett, J.~Ahmed, and A.~Overwijk.
Approximate nearest neighbor negative contrastive learning for dense text retrieval.
In \textit{ICLR}, 2021.

\bibitem{cocohanere2025}
Anonymous.
CoCoHaNeRe: Collaborative contrastive hard negative re-mining for code search.
\textit{ACM TOSEM}, 2025.

\bibitem{jimenez2024swerank}
C.~E. Jim\'enez et~al.
SweRank: Code retrieval for software engineering.
\textit{arXiv preprint}, 2024.

\bibitem{jimenez2024swesmith}
C.~E. Jim\'enez et~al.
SWE-smith: Scaling data for software engineering agents.
\textit{arXiv preprint}, 2024.

\bibitem{guo2021graphcodebert}
D.~Guo et~al.
GraphCodeBERT: Pre-training code representations with data flow.
In \textit{ICLR}, 2021.

\bibitem{karpukhin2020dense}
V.~Karpukhin, B.~Oguz, S.~Min, P.~Lewis, L.~Wu, S.~Edunov, D.~Chen, and W.-t. Yih.
Dense passage retrieval for open-domain question answering.
In \textit{EMNLP}, 2020.

\bibitem{reimers2019sentence}
N.~Reimers and I.~Gurevych.
Sentence-BERT: Sentence embeddings using siamese BERT-networks.
In \textit{EMNLP-IJCNLP}, 2019.

\end{thebibliography}

\end{document}
